
\begin{abstract}
% This paper presents deterministic Riccati observers on $\SO(3)$ for attitude estimation using only scalar measurements, in contrast to conventional methods that rely on directional measurements. A comprehensive analysis establishes sufficient conditions for uniform observability, ensuring local exponential stability of the proposed observers. The key contribution lies in showing that scalar-measurement-based frameworks can substantially enhance the robustness and general applicability of attitude estimation designs. Experimental results validate the effectiveness, versatility, and robustness of the approach across various operating conditions, including scenarios with limited scalar measurements and a biased gyroscope.
Attitude estimation methods typically rely on full vector measurements from inertial sensors such as accelerometers and magnetometers. This paper shows that reliable estimation can also be achieved using only scalar measurements, which naturally arise either as components of vector readings or as independent constraints from other sensing modalities. We propose nonlinear deterministic observers on $\SO(3)$ that incorporate gyroscope bias compensation and guarantee uniform local exponential stability under suitable observability conditions. A key feature of the framework is its robustness to partial sensing: accurate estimation is maintained even when only a subset of vector components is available. Experimental validation on the BROAD dataset confirms consistent performance across progressively reduced measurement configurations, with estimation errors remaining small even under severe information loss. To the best of our knowledge, this is the first work to establish fundamental observability results showing that two scalar measurements under suitable excitation suffice for attitude estimation, and that three are enough in the static case. These results position scalar-measurement-based observers as a practical and reliable alternative to conventional vector-based approaches.


\end{abstract}

\begin{IEEEkeywords}
Uniform Observability, Scalar Measurements, Observers for Nonlinear Systems, Continuous Riccati Equation (CRE).
\end{IEEEkeywords}


\section{Introduction}
\subsection{Motivation}
Precise attitude (orientation) estimation represents a critical capability for a wide range of robotic and aerospace systems, including spacecraft stabilization, aerial and ground vehicle control, and autonomous navigation \cite{crassidis2007survey}.
 In practice, attitude estimation relies on sensors capable of measuring fixed inertial directions in the body frame, typically accelerometers and magnetometers integrated into compact Inertial Measurement Units (IMUs). Accelerometers capture the apparent gravity vector and thus enable roll and pitch estimation under quasi-static conditions, where translational accelerations are negligible. Magnetometers complement this information by observing the Earth’s magnetic field, which provides a reference for heading 
estimation when measurements are not affected by local magnetic interference. These vector measurements, when combined with gyroscopes—also integrated within the IMU—that provide high-rate angular velocity data, constitute the core sensing suite employed in most attitude estimation frameworks. Despite their low cost and wide availability, IMUs are known to suffer from limited accuracy and high sensitivity to noise. In particular, gyroscopes are prone to noise and bias, which can significantly degrade estimation accuracy over time \cite{lefferts1982kalman}. To overcome these limitations, additional sensors are often incorporated. The Global Positioning System (GPS) provides absolute position and velocity references, while Pitot tubes and barometer deliver airspeed and altitude information useful for attitude  estimation of aerial vehicles \cite{Berkane_cdc_2017_gps,tchonkeu2025baromete,DEOLIVEIRA2024_pitot}. Vision-based sensors offer rich geometric constraints that can be exploited for attitude estimation \cite{Hamel_CDC_2020_Vision_aided}. The fusion of these heterogeneous measurements within dedicated estimation frameworks enhances robustness and accuracy by compensating for the individual weaknesses of each sensing modality.
 
This paper revisits the scalar-measurement framework for rigid-body orientation estimation originally introduced in~\cite{alnahhal2025scalar}. In this setting, measurements are typically expressed as the cosine of the angle between two vectors, but they can also arise from alternative constraints such as tilt relationships derived from barometric or range data, or altitude differences inferred from landmarks. Unlike full vector observations, scalar measurements provide only partial constraints about the attitude; nevertheless, they allow for a more flexible and resilient estimation process capable of incorporating measurements from a wide variety of sensors, including those unable to provide complete directional data. This versatility not only enables the integration of innovative sensor combinations but also proves particularly advantageous in scenarios where traditional vector observations are unreliable or incomplete due to sensor degradation, magnetic disturbances, occlusions in vision-based systems, or simplified sensing architectures designed for compact or low-cost platforms, thereby enhancing the overall robustness and accuracy of attitude estimation across diverse applications.

\subsection{Contributions of the Paper}
While the existing formulation in~\cite{alnahhal2025scalar} demonstrated the feasibility of attitude estimation from scalar measurement, its implementation relied on embedding the rotation group $\SO(3)$ into the Euclidean space $\mathbb{R}^9$. This representation introduced redundant states, imposed unnecessarily strict uniform observability conditions, and neglected gyroscope bias compensation. To overcome these limitations, we adopt a direct design on $\SO(3)$ that explicitly incorporates gyroscope bias while preserving the intrinsic geometric structure of rotations. Building on the Riccati-based observer framework of~\cite{hamel2018riccati}, the proposed deterministic filter is inspired by the multiplicative Extended Kalman Filter (MEKF), in which the continuous Riccati equation (CRE) governs both the observer dynamics and the Lyapunov-based stability analysis~\cite{hamel2017position}. 

The main contributions of this work can be summarized as follows. Firstly, we establish a direct $\SO(3)$-based observer that avoids high-dimensional embeddings while ensuring a geometrically consistent attitude representation and explicit gyroscope bias handling. Secondly, rigorous persistence-of-excitation conditions are derived, linking the observability of the linearized dynamics to the uniform local exponential stability of the attitude estimation error. Thirdly, new theoretical insights are provided for scalar-based attitude estimation. In particular, while it is well known that at least two non-collinear vector measurements are required to reconstruct attitude~\cite{Mahony_Hamel_Pflimlin,batista2012sensor}, we show that two scalar measurements under suitable excitation are sufficient to ensure observability, and that three scalar measurements are enough in static conditions, consistent with the intrinsic three-dimensional nature of $\SO(3)$. To the best of our knowledge, these results have not been reported before. The effectiveness and robustness of the proposed approach are further demonstrated through experimental validation under partial and noisy measurement conditions.
%**To achieve accurate attitude estimatoin, reliable state observers (filters) have been extensively studied, aiming to fuse information from multiple sensors—typically providing vector inertial measurements, expressed in either the inertial or body frame, together with body-frame angular velocity measurements \cite{Mahony_Hamel_Pflimlin,batista2012sensor,zlotnik2016nonlinear,tayebi2006attitude,izadi2014rigid}.**
\section{Related Work}
The attitude estimation problem has its roots in the classical attitude determination formulation introduced by Wahba in 1965 \cite{wahba1965least}. Known as Wahba’s problem, it seeks the rotation matrix that best aligns a set of measured body-frame vectors with their corresponding inertial references through a weighted least-squares criterion. 
Over the past decades, this fundamental idea has evolved through successive generations of methods ranging from deterministic reconstruction algorithms to stochastic filters and nonlinear geometric observers. The following subsections briefly review these three families of approaches and outline the conceptual and methodological evolution of attitude estimation, leading to the recent geometric and Riccati-based frameworks that inspire the present work.

\subsection{Deterministic Reconstruction Methods}
Early solutions to Wahba’s problem \cite{wahba1965least} focused on deterministic reconstruction methods, such as Davenport’s $q$-method, Shuster’s QUEST algorithm, and Markley’s SVD-based approach \cite{shuster1981three,markley1988attitude}. These techniques aimed to determine the attitude that best aligns a set of measured body-frame vectors with their known inertial counterparts, formulating the problem as a weighted least-squares optimization that minimizes the discrepancy between the two sets of observations. By exploiting algebraic properties of rotation matrices and quaternions, these methods provided closed-form solutions that were particularly suitable for early spacecraft and satellite applications, where computational resources were limited. Despite their precision in static conditions, deterministic reconstruction methods operate on individual measurement sets and thus neglect temporal correlations, process noise, and sensor dynamics. As a result, they are inherently limited to snapshot attitude determination rather than continuous estimation. These constraints motivated the transition toward recursive and stochastic filtering frameworks capable of propagating uncertainty and bias estimates over time.
%Early solutions to Wahba’s problem \cite{wahba1965least} focused on deterministic reconstruction methods, such as Davenport’s $q$-method, Shuster’s QUEST algorithm, and Markley’s SVD-based approach \cite{shuster1981three,markley1988attitude}. These techniques provide closed-form or iterative solutions to the attitude estimation problem from vector observations, but they do not explicitly handle sensor noise or dynamic models. 
\subsection{Kalman-type Filters}
The Kalman family of filters remains the most widely adopted framework for attitude estimation in aerospace and robotics, owing to its balance between statistical optimality, computational tractability, and practical robustness. The Extended Kalman Filter (EKF) represents the classical approach to nonlinear estimation by locally linearizing system and measurement models around the current estimate \cite{lefferts1982kalman,crassidis2007survey,barrau2017invariant,barrau2018invariant,Unscented_Julier_Uhlmann_1997_Aerosense,Unscented_Julier_Uhlmann_DurrantWhyte_ICRA_2000,Unscented_Crassidis_Markley_JGCD_2003,Unscented_Julier_Uhlmann_ProcIEEE_2004}. In attitude estimation, the EKF fuses high-rate gyroscopic integration with vector observations such as gravity and geomagnetic field directions to correct long-term drift. However, because attitude evolves on the nonlinear manifold $\SO(3)$, a direct Euclidean linearization can lead to constraint violations in rotation representations.

To overcome this geometric inconsistency, MEKF introduces a small-angle error on SO(3) and applies corrections multiplicatively, ensuring proper manifold evolution \cite{lefferts1982kalman}. The MEKF remains the benchmark for spacecraft and aerial vehicle attitude determination, offering improved numerical stability and consistent estimation of both attitude and gyro bias.
Moreover, the Invariant Extended Kalman Filter (IEKF) further exploits the Lie group structure of the attitude kinematics by defining the estimation error directly on the group 
$\SO(3)$, rather than on its local coordinates \cite{barrau2017invariant,barrau2018invariant}. This invariant formulation yields error dynamics that are independent of the chosen attitude parameterization and remain equivariant under changes of reference frame, providing stronger theoretical guarantees of local exponential stability and improved robustness in highly nonlinear or fast-rotating regimes. Besides, to address the limitations of local linearization and enhance performance under strong nonlinearities, the Unscented Kalman Filter (UKF) replaces Jacobian linearization with the propagation of deterministically chosen sigma points through the nonlinear dynamics, capturing mean and covariance up to second-order accuracy \cite{Unscented_Julier_Uhlmann_1997_Aerosense,Unscented_Julier_Uhlmann_DurrantWhyte_ICRA_2000,Unscented_Crassidis_Markley_JGCD_2003,Unscented_Julier_Uhlmann_ProcIEEE_2004}.
%To address these aspects, Kalman-type filters were later introduced for attitude and bias estimation \cite{lefferts1982kalman}. However, these filters are computationally intensive and rely on linear approximations, requiring careful tuning and implementation~\cite{crassidis2007survey}. To overcome such limitations, invariant Kalman filters have emerged as a robust alternative, offering local asymptotic stability and improved performance in nonlinear settings~\cite{barrau2018invariant,bonnabel2009invariant}. 


\subsection{Nonlinear Deterministic Filters} 
In parallel, nonlinear deterministic observers have been developed, marking a shift from statistical optimality toward geometric consistency and stability guarantees. Early work in geometric control established the theoretical foundations for stabilization on Lie groups~\cite{Jurdjevic_1972,Bullo_1999,Morin_2003}. Building on these principles, nonlinear complementary filters on SO(3) were introduced to ensure almost global asymptotic stability ~\cite{Mahony_Hamel_Pflimlin,Baldwin_2007}. 
Following these developments, other works focused on enhancing global convergence properties and addressing gyroscope bias estimation~\cite{tayebi2006attitude,Sanyal_2008,Lageman_2009,Vasconcelos_2010,hua2014implementation,batista2012sensor,zlotnik2016nonlinear}. More recently, attention has shifted toward handling time-varying reference vectors and establishing rigorous links between uniform observability and estimator stability~\cite{grip2011attitude,Trumpf_2012}. This analysis underpins the Riccati-based observer framework, which utilizes the Continuous Riccati Equation (CRE) to address nonstationary problems such as the Perspective-n-Point (PnP) problem~\cite{hamel2018riccati}. This geometric framework has also been extended to vision-aided INS, fusing inertial measurements with visual bearing-to-landmark measurements to robustly estimate the full navigation state~\cite{Wang_2022}.

While these approaches have shown promising results, they typically assume the availability of complete three-dimensional vector measurements—an assumption that may not hold in practice due to sensor noise, environmental disturbances, or hardware limitations.

